\subsection*{一步一步建立直观}

\paragraph{1. 从屏幕图像恢复空间对象}

旋转观察视角，留意箭头在屏幕上变短并不意味着空间长度变小。屏幕只是三维对象的投影，下面的长度读数才使用三维坐标计算。操作参数改变数学对象，拖动空白处改变相机，这两件事必须明确区分。



\paragraph{2. 绕固定轴旋转改变哪些坐标}

绕 $z$ 轴转四分之一圈时，$(x,y,z)$ 变为 $(-y,x,z)$。横纵分量互换并改变一个符号，高度保持。平方长度仍是 $x^2+y^2+z^2$，因此图像中的转动有一条可以独立核对的代数不变量。可以在任意角度先记录长度，再修改转角复核这个守恒关系。



\paragraph{3. 旋转与投影的本质差别}

切换到投影模式，垂直虚线显示被丢弃的高度。所有只在 $z$ 坐标上不同的点都有同一个平面像。旋转保留恢复原状态的信息，投影则一般做不到；能否求逆取决于这个信息差别，而不是画面是否仍有三维感。
